\chapter{کار با متن}
\label{ch:strings}

\begin{goals}
\begin{itemize}
  \item چسباندن رشته‌ها و ساختن پیام‌های ترکیبی
  \item نویسه‌های ویژه: خط جدید، جدول‌بندی و نقل قول درون رشته
  \item بریدن بخشی از رشته، شکافتن به واژه‌ها و پیراستن فاصله‌ها
  \item جست‌وجو در رشته و مقایسه‌ی رشته‌ها
  \item تبدیل رشته به عدد و برعکس
\end{itemize}
\end{goals}

\section{متن همه‌جا هست}

نام‌ها، پیام‌ها، نشانی‌ها، شماره‌تلفن‌ها، تاریخ‌ها: بخش بزرگی از داده‌های
هر برنامه متن است. در فصل‌های پیش با رشته‌ها کار کرده‌ایم، اما همیشه به
عنوان یک تکه‌ی یکپارچه. در این فصل یاد می‌گیریم درون رشته را ببینیم، از آن
تکه‌ای برداریم، آن را بشکافیم و از آن اطلاعات بیرون بکشیم. همه‌ی این کارها
با \textbf{قابلیت‌هایی} انجام می‌شوند که با یک نقطه به رشته می‌چسبند، مثل
\fcode{به صحیح} که در فصل \ref{ch:types} دیدید.

\section{چسباندن و طول}

\program{چسباندن رشته‌ها}
\begin{salam}
تابع اصلی:
    نام := "سارا"
    نام خانوادگی := "احمدی"
    نام کامل := نام + " " + نام خانوادگی
    چاپ نام کامل
    شماره := "09121234567"
    چاپ "طول شماره:", شماره.طول()
تمام
\end{salam}

\begin{outputfa}
سارا احمدی
طول شماره: 11
\end{outputfa}

\fcode{طول()} تعداد نویسه‌های رشته را می‌دهد. برای شماره‌تلفن یا هر متنی
که از رقم و حرف انگلیسی ساخته شده، این عدد همان چیزی است که انتظار دارید.

\begin{warn}[طول متن فارسی]
رایانه هر حرف فارسی را در \emph{دو} خانه نگه می‌دارد، در حالی که هر رقم،
حرف انگلیسی یا فاصله یک خانه می‌گیرد. \fcode{طول()} تعداد خانه‌ها را
می‌شمارد؛ پس \fcode{"سلام".طول()} برابر ۸ است، نه ۴. به همین دلیل در
این کتاب \fcode{طول()} و \fcode{زیررشته()} را بیشتر روی متن‌هایی به کار
می‌بریم که از رقم و حرف انگلیسی ساخته شده‌اند (شماره، تاریخ، کد). سلام
برای شمردن دقیق حرف‌های فارسی ابزارهای دیگری دارد که به این کتاب مربوط
نیست.
\end{warn}

\section{نویسه‌های ویژه}

بعضی چیزها را نمی‌توان مستقیم درون نقل قول نوشت: رفتن به خط بعد، یا خود
علامت نقل قول. برای این‌ها سلام ترکیب‌هایی با خط مورب وارونه \code{\textbackslash}
دارد:

\begin{center}
\begin{tabular}{cl}
\toprule
نوشتار & معنی \\
\midrule
\code{\textbackslash n} & رفتن به خط بعد \\
\code{\textbackslash t} & جدول‌بندی (یک پرش بلند) \\
\code{\textbackslash "} & خود علامت نقل قول \\
\code{\textbackslash\textbackslash} & خود خط مورب وارونه \\
\bottomrule
\end{tabular}
\end{center}

\program{نویسه‌های ویژه}
\begin{salam}
تابع اصلی:
    چاپ "خط اول\nخط دوم"
    چاپ "نام:\tسارا"
    چاپ "او گفت: \"سلام\""
تمام
\end{salam}

\begin{outputfa}
خط اول
خط دوم
نام:	سارا
او گفت: "سلام"
\end{outputfa}

با \code{\textbackslash n} یک \fcode{چاپ} می‌تواند چند خط بنویسد. این
ترکیب‌ها هر کدام \emph{یک} نویسه به حساب می‌آیند، هرچند با دو علامت نوشته
می‌شوند.

\section{بریدن بخشی از رشته}

\fcode{زیررشته(آغاز, تعداد)} از جایگاه \fcode{آغاز} (که مثل آرایه‌ها از صفر
شمرده می‌شود) به تعداد \fcode{تعداد} نویسه برمی‌دارد:

\program{تکه‌های یک تاریخ}
\begin{salam}
تابع اصلی:
    تاریخ := "1404/06/28"
    سال := تاریخ.زیررشته(0, 4)
    ماه := تاریخ.زیررشته(5, 2)
    روز := تاریخ.زیررشته(8, 2)
    چاپ "سال:", سال, "ماه:", ماه, "روز:", روز
    چاپ "سال بعد:", سال.به صحیح() + 1
تمام
\end{salam}

\begin{outputfa}
سال: 1404 ماه: 06 روز: 28
سال بعد: 1405
\end{outputfa}

جایگاه‌ها را بشمارید: چهار رقم سال در جایگاه‌های ۰ تا ۳ هستند، بعد یک
خط مورب در جایگاه ۴، بعد ماه در ۵ و ۶، خط مورب در ۷ و روز در ۸ و ۹.
نتیجه‌ی \fcode{زیررشته} یک رشته‌ی تازه است؛ رشته‌ی اصلی دست‌نخورده می‌ماند.

با همین ابزار می‌توان یک شماره‌تلفن را بررسی کرد:

\program{بررسی شماره‌ی همراه}
\begin{salam}
تابع اصلی:
    شماره := "09121234567"
    پیش شماره := شماره.زیررشته(0, 4)
    اگر شماره.طول() == 11 و پیش شماره == "0912":
        چاپ "شماره‌ی همراه اول است"
    وگرنه:
        چاپ "شماره‌ی ناشناخته"
    تمام
تمام
\end{salam}

\begin{outputfa}
شماره‌ی همراه اول است
\end{outputfa}

\section{شکافتن به واژه‌ها}

\fcode{بشکاف(جداکننده)} رشته را در هر جایی که جداکننده هست می‌برد و
فهرستی از تکه‌ها می‌دهد. این فهرست شبیه آرایه است، با این تفاوت که عضوها
را با \fcode{بگیر(اندیس)} می‌خوانیم:

\program{شمارش واژه‌های یک جمله}
\begin{salam}
تابع اصلی:
    جمله := "زبان سلام برای همه است"
    کلمات := جمله.بشکاف(" ")
    چاپ "تعداد واژه‌ها:", کلمات.طول()
    تکرار کلمات.طول() با شماره:
        چاپ شماره + 1, "-", کلمات.بگیر(شماره)
    تمام
تمام
\end{salam}

\begin{outputfa}
تعداد واژه‌ها: 5
1 - زبان
2 - سلام
3 - برای
4 - همه
5 - است
\end{outputfa}

شکافتن با فاصله روی متن فارسی هم درست کار می‌کند، چون فاصله همیشه یک
نویسه است. جداکننده می‌تواند هر رشته‌ای باشد: \fcode{بشکاف("،")} برای
فهرستی که با ویرگول فارسی جدا شده، یا \fcode{بشکاف("/")} برای همان تاریخ
بالا.

\section{پیراستن فاصله‌های اضافی}

متنی که از کاربر یا از یک پرونده می‌آید، اغلب فاصله‌های اضافی در آغاز و
پایان دارد. \fcode{پیراست()} آن‌ها را می‌زداید:

\program{پیراستن}
\begin{salam}
تابع اصلی:
    ورودی := "   قبول   "
    پاک := ورودی.پیراست()
    چاپ "[" + ورودی + "]"
    چاپ "[" + پاک + "]"
    چاپ پاک == "قبول"
تمام
\end{salam}

\begin{outputfa}
[   قبول   ]
[قبول]
true
\end{outputfa}

کروشه‌ها را برای این گذاشتیم که فاصله‌های نامرئی دیده شوند. بدون
پیراستن، مقایسه با \code{"قبول"} نادرست می‌شد، چون فاصله‌ها هم بخشی از
رشته‌اند. قاعده‌ی خوب: هر متنی را که از بیرون می‌آید، پیش از مقایسه
بپیرایید.

\section{جست‌وجو در رشته}

\fcode{بیاب(متن)} می‌گوید متن داده‌شده در کجای رشته شروع می‌شود، و اگر
اصلاً در رشته نباشد، \code{-1} برمی‌گرداند. رایج‌ترین کاربردش همین پرسش
«آیا هست یا نه» است:

\program{جست‌وجو در پیام}
\begin{salam}
تابع اصلی:
    پیام := "پرداخت با موفقیت انجام شد"
    اگر پیام.بیاب("موفقیت") != -1:
        چاپ "عملیات موفق بود"
    وگرنه:
        چاپ "عملیات ناموفق بود"
    تمام
    چاپ پیام.بیاب("خطا")
تمام
\end{salam}

\begin{outputfa}
عملیات موفق بود
-1
\end{outputfa}

\section{رشته و عدد}

در فصل \ref{ch:types} دیدیم که \code{+} میان رشته و عدد، عدد را به رشته
تبدیل می‌کند و می‌چسباند، و \fcode{به صحیح()} و \fcode{به اعشار()} راه
عکس را می‌روند:

\program{از متن به عدد}
\begin{salam}
تابع اصلی:
    قیمت متن := "25000"
    تعداد متن := "3"
    قیمت := قیمت متن.به صحیح()
    تعداد := تعداد متن.به صحیح()
    چاپ "مبلغ:", قیمت * تعداد
    وزن := "2.5".به اعشار()
    چاپ "وزن دو کارتن:", وزن * 2.0
تمام
\end{salam}

\begin{outputfa}
مبلغ: 75000
وزن دو کارتن: 5
\end{outputfa}

\begin{warn}
\fcode{به صحیح()} فقط رقم‌های انگلیسی را می‌فهمد. رشته‌ای مثل
\code{"۲۵۰۰۰"} با رقم‌های فارسی به صفر تبدیل می‌شود، نه به ۲۵٬۰۰۰. اگر
متنی با رقم فارسی دارید، پیش از تبدیل باید رقم‌هایش را انگلیسی کنید؛
سلام برای این کار ابزار آماده دارد که در کتاب‌های بعدی می‌بینید.
\end{warn}

\section{ساختن رشته با حلقه}

در فصل \ref{ch:loops} چند بار رشته‌ای را قطعه‌قطعه ساختیم. این کار را
می‌توان در تابعی جمع کرد که هر رشته‌ای را هر تعداد بار تکرار کند:

\program{تکرار یک رشته}
\begin{salam}
تابع چندبار(متن: رشته, تعداد: صحیح): رشته:
    متغیر نتیجه := ""
    تکرار تعداد:
        نتیجه += متن
    تمام
    بازگشت نتیجه
تمام

تابع اصلی:
    چاپ چندبار("-=", 10)
    چاپ چندبار("*", 5) + " پنج ستاره " + چندبار("*", 5)
    چاپ چندبار("-=", 10)
تمام
\end{salam}

\begin{outputfa}
-=-=-=-=-=-=-=-=-=-=
***** پنج ستاره *****
-=-=-=-=-=-=-=-=-=-=
\end{outputfa}

\section{یک برنامه‌ی کامل‌تر}

تاریخی به شکل \code{"1404/06/28"} را می‌گیریم و آن را به شکل خوانا، با
نام ماه، چاپ می‌کنیم:

\program{تاریخ خوانا}
\begin{salam}
تابع نام ماه(شماره: صحیح): رشته:
    بازگشت تطبیق شماره:
        1 => "فروردین"
        2 => "اردیبهشت"
        3 => "خرداد"
        4 => "تیر"
        5 => "مرداد"
        6 => "شهریور"
        7 => "مهر"
        8 => "آبان"
        9 => "آذر"
        10 => "دی"
        11 => "بهمن"
        12 => "اسفند"
        وگرنه => "نامعتبر"
    تمام
تمام

تابع اصلی:
    تاریخ := "1404/06/28"
    سال := تاریخ.زیررشته(0, 4)
    ماه := تاریخ.زیررشته(5, 2).به صحیح()
    روز := تاریخ.زیررشته(8, 2).به صحیح()
    چاپ روز, نام ماه(ماه), سال
تمام
\end{salam}

\begin{outputfa}
28 شهریور 1404
\end{outputfa}

به خط \fcode{تاریخ.زیررشته(5, 2).به صحیح()} دقت کنید: نتیجه‌ی
\fcode{زیررشته} یک رشته است و می‌توان بلافاصله قابلیت دیگری را روی آن
صدا زد. این «زنجیره کردن» در کار با رشته‌ها بسیار رایج است.

\begin{summary}
\begin{itemize}
  \item \code{+} رشته‌ها را می‌چسباند و \fcode{طول()} تعداد نویسه‌ها را
        می‌دهد (هر حرف فارسی دو تا شمرده می‌شود).
  \item \code{\textbackslash n}، \code{\textbackslash t} و
        \code{\textbackslash "} نویسه‌های ویژه‌اند.
  \item \fcode{زیررشته(آغاز, تعداد)} تکه‌ای را می‌برد،
        \fcode{بشکاف(جداکننده)} رشته را به فهرستی از تکه‌ها می‌شکافد،
        \fcode{پیراست()} فاصله‌های دو سر را می‌زداید و \fcode{بیاب(متن)}
        جایگاه یک متن را (یا \code{-1}) می‌دهد.
  \item \fcode{به صحیح()} و \fcode{به اعشار()} رشته را به عدد تبدیل
        می‌کنند (فقط با رقم‌های انگلیسی).
  \item قابلیت‌ها را می‌توان زنجیره کرد: \fcode{متن.زیررشته(0, 4).به صحیح()}.
\end{itemize}
\end{summary}

\section*{تمرین‌ها}

\exercise نام و نام خانوادگی را در دو متغیر بگذارید و با یک \fcode{چاپ}
سه خط چاپ کنید: نام کامل، نام خانوادگی به تنهایی، و جمله‌ای که هر دو را
در خود دارد. (راهنمایی: \code{\textbackslash n}.)

\exercise کد ملی ده‌رقمی را در رشته‌ای بگذارید و بررسی کنید که دقیقاً ده
نویسه دارد. سپس سه رقم اول آن را جدا کنید و چاپ کنید.

\exercise جمله‌ای را در متغیری بگذارید و درازترین واژه‌ی آن را پیدا و
چاپ کنید. (چون واژه‌ها فارسی‌اند، \fcode{طول()} هر واژه دو برابر شمار
حرف‌های آن است؛ برای مقایسه‌ی درازی این مشکلی ایجاد نمی‌کند.)

\exercise فهرستی از نمره‌ها را به شکل یک رشته‌ی جداشده با ویرگول
(\code{"17,15,19,12"}) در متغیری بگذارید. آن را بشکافید، هر تکه را به
عدد تبدیل کنید و مجموع و میانگین را چاپ کنید.

\exercise تابعی بنویسید که یک رشته و یک عرض بگیرد و اگر رشته کوتاه‌تر از
عرض بود، به سمت راست آن آن‌قدر فاصله بچسباند که به آن عرض برسد. با آن
یک جدول دوستونی مرتب از نام کالا و قیمت چاپ کنید. (از رشته‌های انگلیسی
یا رقم استفاده کنید تا \fcode{طول()} دقیق باشد.)
